\section{Analysis of Model~$B$ ($\gamma_1,\gamma_2>0$, $\theta_1=\theta_2=0$)}\label{sec:model_b}

Model~$B$ introduces bidirectional jockeying ($\gamma_1,\gamma_2>0$) while retaining no abandonment ($\theta_1=\theta_2=0$). Jockeying conserves the total number of customers in the system, so the stationary distribution still exists under the same condition as Model~$A$. However, jockeying breaks the per-class Poisson/renewal structure: a Class-2 customer can no longer identify its effective service time with a Class-1 busy period, so the Pollaczek--Khinchine probabilistic route is unavailable. Formally integrating the fundamental equation by the method of characteristics yields a general solution, but closing it for $P_y(y)$ requires a Volterra-type integral equation that does not admit a closed form. Model~$B$ is therefore intentionally left open, in structural parallel with Model~$X$; its role is to motivate the tractable specialization Model~$B_2$.

\textbf{Stability.} $\rho = (\lambda_1+\lambda_2)/\mu < 1$. Jockeying does not affect this condition.

\begin{figure}[H]
    \centering
    \resizebox{0.62\textwidth}{!}{%
        % Model B: θ₁=θ₂=0, γ₁,γ₂>0
% Bidirectional jockeying between queues (dashed). No abandonment.
\begin{tikzpicture}[>=Latex, thick]

    % ── Queue 1 (top, Class-1, priority) ──────────────────────────────────────
    \draw (0, 1.9) -- (2, 1.9) -- (2, 1.1) -- (0, 1.1);
    \foreach \x in {0.2, 0.4, ..., 1.8} { \draw[thin] (\x, 1.8) -- (\x, 1.2); }
    \draw[->] (-1.5, 1.5) -- (-0.1, 1.5) node[midway, above] {$\lambda_1$};

    % ── Queue 2 (bottom, Class-2) ──────────────────────────────────────────────
    \draw (0, -1.1) -- (2, -1.1) -- (2, -1.9) -- (0, -1.9);
    \foreach \x in {0.2, 0.4, ..., 1.8} { \draw[thin] (\x, -1.2) -- (\x, -1.8); }
    \draw[->] (-1.5, -1.5) -- (-0.1, -1.5) node[midway, above] {$\lambda_2$};

    % ── Jockeying between queues ──────────────────────────────────────────────
    % γ₁: Class-1 jockeys down to Class-2 queue
    \draw[->, dashed] (0.6, 1.0) -- (0.6, -1.0) node[midway, left] {$\gamma_1$};
    % γ₂: Class-2 jockeys up to Class-1 queue
    \draw[->, dashed] (1.4, -1.0) -- (1.4, 1.0) node[midway, right] {$\gamma_2$};

    % ── Server ────────────────────────────────────────────────────────────────
    \node[draw, circle, minimum size=1.2cm] (server) at (5, 0) {$\mu$};
    \draw[->] (2.2,  1.5) -- (server);
    \draw[->] (2.2, -1.5) -- (server);
    \draw[->] (server) -- (7, 0);

\end{tikzpicture}
%
    }
    \caption{Queue layout for Model~$B$. Dashed arrows indicate bidirectional
    jockeying: $\gamma_1$ moves a Class-1 customer down to Queue~2 and $\gamma_2$
    moves a Class-2 customer up to Queue~1. Both transfers conserve the total~$n$.}
    \label{fig:model-b-queue}
\end{figure}

\begin{claim}\label{clm:B:volterra}
    Under stability $\rho<1$, the bivariate PGF of Model~$B$ satisfies, on each characteristic curve $\gamma_2 x + \gamma_1 y = C_1$ of the fundamental equation, the integro-differential representation
    \begin{equation}
        P(x,y(x)) \;=\; P_h(x,y(x))\int_{0}^{x}
        \frac{h\bigl(t;\,C_1\bigr)}{P_h\bigl(t,\,y(t)\bigr)}\,dt,
        \label{eq:B:integral_rep}
    \end{equation}
    where $P_h$ is the explicit homogeneous solution derived by the method of characteristics~\eqref{eq:B:Ph}, $y(t)=(C_1-\gamma_2 t)/\gamma_1$ is the characteristic parametrisation, and $h(t;C_1)$ is a forcing function depending on the unknown boundary trace $P_y(y)=P(0,y)$.  Evaluating~\eqref{eq:B:integral_rep} at the diagonal $x=y=z$ and invoking the conservation identity $P(z,z)=\pi(0,0)/(1-\rho z)$ reduces the problem to a \emph{Volterra integral equation of the first kind} for $P_y$:
    \begin{equation}
        \mathcal{F}(z) \;=\; \int_0^z \mathcal{K}(z,t)\,P_y\!\bigl(y(t;z)\bigr)\,dt,
        \label{eq:B:volterra}
    \end{equation}
    where $\mathcal{K}$ and $\mathcal{F}$ are fully explicit (defined in~\eqref{eq:B:kernel}--\eqref{eq:B:rhs} below).  The kernel $\mathcal{K}$ does not factorise for general $(\gamma_1,\gamma_2)$, so~\eqref{eq:B:volterra} does not reduce to a closed-form solution for $P_y$.  Model~$B$ is therefore left open.
\end{claim}

\subsection{Analysis of Model~$B$ by the Method of Characteristics and Variation of Parameters}

\subsubsection*{Homogeneous solution via the method of characteristics}

The fundamental equation~\eqref{eq:gen:fundamental_equation} with $\theta_1=\theta_2=0$ is a first-order linear PDE in $P(x,y)$.  Setting the right-hand side to zero and dividing through by $y$ yields the homogeneous PDE
\begin{equation}
    \gamma_1 x(x-y)\,\frac{\partial P}{\partial x} - \gamma_2 x(x-y)\,\frac{\partial P}{\partial y}
    = -\bigl[(\lambda_1+\lambda_2+\mu)x - \mu - \lambda_1 x^2 - \lambda_2 xy\bigr]P.
    \label{eq:B:homogeneous_pde}
\end{equation}
The method of characteristics converts this into the system of ODEs
\begin{equation}
    \frac{dx}{\gamma_1 x(x-y)}
    = \frac{dy}{-\gamma_2 x(x-y)}
    = \frac{dP}{-\bigl[(\lambda_1+\lambda_2+\mu)x - \mu - \lambda_1 x^2 - \lambda_2 xy\bigr]P}.
    \label{eq:B:char_system}
\end{equation}
Taking the first ratio, the $x(x-y)$ factors cancel:
\begin{equation}
    \frac{dx}{\gamma_1} = \frac{dy}{-\gamma_2}
    \quad\Longrightarrow\quad
    C_1 = \gamma_2 x + \gamma_1 y \;\text{ (constant along each characteristic)}.
    \label{eq:B:char_constant}
\end{equation}
Parameterising $y$ along the curve by $y(x) = (C_1 - \gamma_2 x)/\gamma_1$ and substituting into the $dP/P$ ratio:
\begin{align}
    \frac{dP}{P}
    &= \frac{(\gamma_1\lambda_1 - \lambda_2\gamma_2)x^2
            - \bigl[\gamma_1(\lambda_1+\lambda_2+\mu) - \lambda_2 C_1\bigr]x
            + \gamma_1\mu}
           {\gamma_1 x\bigl[(\gamma_1+\gamma_2)x - C_1\bigr]}\,dx.
    \label{eq:B:dPP}
\end{align}
Partial fraction decomposition of the right-hand side, writing the integrand as $(1/\gamma_1)\,I(x)\,dx$ with
\begin{equation}
    I(x) = K + \frac{M}{x} + \frac{N(C_1)}{(\gamma_1+\gamma_2)x - C_1},
\end{equation}
yields, by matching coefficients,
\begin{equation}
    K = \frac{\gamma_1\lambda_1 - \lambda_2\gamma_2}{\gamma_1+\gamma_2},
    \qquad
    M = -\frac{\gamma_1\mu}{C_1},
    \qquad
    N(C_1) = \gamma_1\!\left[
        \frac{\lambda_1+\lambda_2}{\gamma_1+\gamma_2}C_1
        - (\lambda_1+\lambda_2+\mu)
        + \frac{\mu(\gamma_1+\gamma_2)}{C_1}
    \right].
    \label{eq:B:KMN}
\end{equation}
Integrating term by term gives
\begin{equation}
    \ln P = \frac{K}{\gamma_1}x - \frac{\mu}{C_1}\ln x
    + \frac{N(C_1)}{\gamma_1(\gamma_1+\gamma_2)}\ln\bigl[(\gamma_1+\gamma_2)x - C_1\bigr]
    + \ln f(C_1),
    \label{eq:B:lnP}
\end{equation}
where $f(C_1)$ is an arbitrary function of the characteristic constant.  Now observe that
\begin{equation}
    (\gamma_1+\gamma_2)x - C_1
    = (\gamma_1+\gamma_2)x - (\gamma_2 x + \gamma_1 y)
    = \gamma_1(x-y),
    \label{eq:B:simplify}
\end{equation}
so $\ln[(\gamma_1+\gamma_2)x-C_1] = \ln\gamma_1 + \ln(x-y)$.  The constant $\ln\gamma_1$ is absorbed into $f(C_1)$ to form a new arbitrary function $F(C_1)$.  Reading off the exponent of $(x-y)$ from~\eqref{eq:B:lnP}:
\begin{equation}
    E(x,y) := \frac{N(C_1)}{\gamma_1(\gamma_1+\gamma_2)}\bigg|_{C_1=\gamma_2 x+\gamma_1 y}
    = \frac{\lambda_1+\lambda_2}{(\gamma_1+\gamma_2)^2}(\gamma_2 x+\gamma_1 y)
    - \frac{\lambda_1+\lambda_2+\mu}{\gamma_1+\gamma_2}
    + \frac{\mu}{\gamma_2 x+\gamma_1 y}.
    \label{eq:B:E}
\end{equation}
Substituting $C_1=\gamma_2 x+\gamma_1 y$ and taking the exponential of~\eqref{eq:B:lnP} yields the general solution to the homogeneous PDE:
\begin{equation}
    P_h(x,y) = F(\gamma_2 x+\gamma_1 y)
    \cdot \exp\!\left(\frac{\gamma_1\lambda_1-\lambda_2\gamma_2}{\gamma_1(\gamma_1+\gamma_2)}x\right)
    \cdot x^{-\mu/(\gamma_2 x+\gamma_1 y)}
    \cdot (x-y)^{E(x,y)},
    \label{eq:B:Ph}
\end{equation}
where $F$ is an arbitrary differentiable function to be determined by the boundary conditions.

\smallskip\noindent\textit{Sign of $E(x,y)$ and connection to Model~$B_2$.}
Under stability ($\rho<1$), a computation analogous to that for $\beta(y)<0$ in Model~$B_2$ shows that $E(x,y)<0$ in the interior of the unit square.  This means $(x-y)^{E(x,y)}\to\infty$ as $x\to y$, which is the analogue of the branch singularity that drives the analyticity argument in Model~$B_2$.  Setting $\gamma_2=0$ in~\eqref{eq:B:Ph} gives $C_1=\gamma_1 y$ (constant in $x$), and $P_h|_{\gamma_2=0}$ reduces to
\begin{equation}
    P_h(x,y)\big|_{\gamma_2=0}
    = F(\gamma_1 y)\cdot e^{\lambda_1 x/\gamma_1}\cdot x^{-\mu/(\gamma_1 y)}\cdot(x-y)^{-\beta(y)},
    \label{eq:B:Ph_B2limit}
\end{equation}
which is $F(\gamma_1 y)\cdot\mu_I^{-1}(x;y)$, the inverse of the Model~$B_2$ integrating factor~\eqref{eq:B2:mu_I}.  The two analyses are therefore structurally identical at $\gamma_2=0$: what is called the "homogeneous solution" in Model~$B$ degenerates to the "inverse integrating factor" in Model~$B_2$.

\subsubsection*{Particular solution via variation of parameters}

Variation of parameters exploits the homogeneous solution as a fundamental solution for the inhomogeneous equation.  Restricting attention to a fixed characteristic curve $y(x)=(C_1-\gamma_2 x)/\gamma_1$, the full fundamental equation~\eqref{eq:gen:fundamental_equation} becomes a first-order ODE in $x$:
\begin{equation}
    \frac{dP}{dx} = f(x;C_1)\,P + h(x;C_1),
    \label{eq:B:ode_char}
\end{equation}
where the coefficient and forcing term are
\begin{align}
    f(x;C_1) &= -\frac{A(x,\,y(x))}{\gamma_1 x\bigl(x-y(x)\bigr)},
    \label{eq:B:fdef}\\[4pt]
    h(x;C_1) &= \frac{G(x,\,y(x))}{\gamma_1 x\,y(x)\bigl(x-y(x)\bigr)},
    \label{eq:B:hdef}
\end{align}
with
\begin{align}
    A(x,y) &= (\lambda_1+\lambda_2+\mu)x - \mu - \lambda_1 x^2 - \lambda_2 xy,
    \label{eq:B:Adef}\\
    G(x,y) &= \mu(x-y)\,P_y(y) - \mu x(1-y)\,\pi(0,0).
    \label{eq:B:Gdef}
\end{align}
Note that $f(x;C_1) = -p(x;y(x))/\gamma_1$ and $h(x;C_1) = q(x;y(x))/y(x)$ in the notation of Model~$B_2$'s ODE~\eqref{eq:B2:pxy}--\eqref{eq:B2:qxy} (evaluated along the characteristic).  Since $P_h(x,y(x))$ satisfies the homogeneous equation $\frac{dP_h}{dx}=f(x;C_1)P_h$, it is the fundamental solution of~\eqref{eq:B:ode_char}.  The general solution of the full inhomogeneous equation is therefore
\begin{equation}
    P(x,y(x)) = P_h(x,y(x))\left[F(C_1) + \int_{x_0}^{x} \frac{h(t;C_1)}{P_h(t,y(t))}\,dt\right].
    \label{eq:B:variation}
\end{equation}

\subsubsection*{Determination of $F(C_1)$ from the boundary condition at $x=0$}

Along any characteristic, $y(0)=C_1/\gamma_1$, so the PGF value at $x=0$ is
\[
    P\bigl(0,\,y(0)\bigr)
    = P\!\left(0,\,\frac{C_1}{\gamma_1}\right)
    = P_y\!\left(\frac{C_1}{\gamma_1}\right),
\]
which must be finite (it is a value of the boundary PGF $P_y$).  However, $P_h$ contains the factor $x^{-\mu/C_1}\to\infty$ as $x\to0^+$ (since $\mu/C_1>0$), while the integral in~\eqref{eq:B:variation} tends to zero (the limits coincide at $x_0=0$).  Hence
\[
    \lim_{x\to0^+}P(x,y(x))
    = \lim_{x\to0^+}P_h(x,y(x))\cdot F(C_1)
    = \infty\cdot F(C_1),
\]
which is finite only if $F(C_1)=0$.  Setting $x_0=0$ and $F(C_1)=0$ in~\eqref{eq:B:variation} gives the representation stated in Claim~\ref{clm:B:volterra}:
\begin{equation}
    P(x,y(x)) = P_h(x,y(x))\int_{0}^{x} \frac{h(t;C_1)}{P_h(t,y(t))}\,dt.
    \label{eq:B:integral_rep_body}
\end{equation}

\subsubsection*{Simplification of the forcing function and diagonal evaluation}

When restricted to the characteristic, $G(t,y(t))$ in~\eqref{eq:B:Gdef} simplifies: the factor $(t-y(t))$ cancels in the first term, and the sign of the second term changes because $y(t)>t$ for $t\in[0,z)$ (so $(t-y(t))<0$).  Explicitly,
\begin{equation}
    h(t;C_1) = \frac{\mu}{\gamma_1\,t\,y(t)}\,P_y\!\bigl(y(t)\bigr)
    + \frac{\mu\,(1-y(t))}{\gamma_1\,y(t)\,(y(t)-t)}\,\pi(0,0).
    \label{eq:B:h_simplified}
\end{equation}

\noindent\textit{Empty-state probabilities.}
Before evaluating at the diagonal, we fix $\pi_0$ and $\pi(0,0)$.  Setting $x=y$ in the full fundamental equation, the factors $(x-y)$ kill both the jockeying convection terms and the boundary term $\mu(x-y)P_y(y)$, leaving
\begin{equation}
    \bigl[(\lambda_1+\lambda_2+\mu)x - \mu - (\lambda_1+\lambda_2)x^2\bigr]P(x,x)
    = -\mu(1-x)\,\pi(0,0).
    \label{eq:B:diagonal_ode}
\end{equation}
The bracket factors as $(\mu-\lambda x)(x-1)$; after cancelling $(x-1)$,
\begin{equation}
    P(x,x) = \frac{\mu\,\pi(0,0)}{\mu-\lambda x} = \frac{\pi(0,0)}{1-\rho x}.
    \label{eq:B:Pxx}
\end{equation}
Setting $x=1$ and using $P(1,1)=1-\pi_0$ together with the idle-state balance $\lambda\pi_0=\mu\pi(0,0)$:
\begin{equation}
    \pi_0 = 1-\rho, \qquad \pi(0,0) = \rho(1-\rho).
    \label{eq:B:pi00}
\end{equation}
This recovers the same values as Model~$A$ and Model~$B_2$~\eqref{eq:B2:pi00}: jockeying conserves the total-customer process and therefore cannot change the marginal idle probability.

\noindent\textit{Diagonal evaluation.}
Every diagonal point $(z,z)$ lies on the characteristic with $C_1=(\gamma_1+\gamma_2)z$.  The parametrisation along that curve is $y(t;z)=((\gamma_1+\gamma_2)z-\gamma_2 t)/\gamma_1$, which satisfies $y(0;z)>z$ and $y(z;z)=z$, confirming that the curve reaches the diagonal.  Evaluating~\eqref{eq:B:integral_rep_body} at $x=z$ and substituting~\eqref{eq:B:Pxx}:
\begin{equation}
    \frac{\pi(0,0)}{1-\rho z}
    = P_h(z,z)\int_0^z \frac{h\bigl(t;\,(\gamma_1+\gamma_2)z\bigr)}{P_h\bigl(t,\,y(t;z)\bigr)}\,dt.
    \label{eq:B:diag_eval}
\end{equation}

\subsubsection*{Reduction to the Volterra equation}

Substituting~\eqref{eq:B:h_simplified} into~\eqref{eq:B:diag_eval} and separating the known term (involving $\pi(0,0)$) from the unknown term (involving $P_y(y(t;z))$):
\begin{align}
    \frac{\pi(0,0)}{1-\rho z}
    &= \frac{\mu\,P_h(z,z)}{\gamma_1}
       \int_0^z \frac{P_y\bigl(y(t;z)\bigr)}{t\cdot y(t;z)\cdot P_h\bigl(t,y(t;z)\bigr)}\,dt
    \nonumber\\
    &\quad
    + \frac{\mu\,P_h(z,z)\,\pi(0,0)}{\gamma_1}
       \int_0^z \frac{1-y(t;z)}{y(t;z)\cdot(y(t;z)-t)\cdot P_h\bigl(t,y(t;z)\bigr)}\,dt.
    \label{eq:B:split}
\end{align}
Define the kernel and the known right-hand side:
\begin{align}
    \mathcal{K}(z,t) &:= \frac{\mu\cdot P_h(z,z)}
                                {\gamma_1\cdot t\cdot y(t;z)\cdot P_h\bigl(t,y(t;z)\bigr)},
    \label{eq:B:kernel}\\[6pt]
    \mathcal{B}(z) &:= \frac{\mu\cdot P_h(z,z)\cdot\pi(0,0)}{\gamma_1}
                        \int_0^z
                        \frac{1-y(t;z)}{y(t;z)\cdot(y(t;z)-t)\cdot P_h\bigl(t,y(t;z)\bigr)}
                        \,dt,
    \label{eq:B:Bdef}
\end{align}
and set $\mathcal{F}(z)=\pi(0,0)/(1-\rho z)-\mathcal{B}(z)$.  Then~\eqref{eq:B:split} becomes
\begin{equation}
    \mathcal{F}(z) = \int_0^z \mathcal{K}(z,t)\,P_y\!\bigl(y(t;z)\bigr)\,dt,
    \label{eq:B:rhs}
\end{equation}
confirming Claim~\ref{clm:B:volterra}.  This is a \emph{Volterra integral equation of the first kind} for the unknown boundary function $P_y$: after the substitution $s=y(t;z)$, the equation takes the standard convolution form, but the kernel $\mathcal{K}$ depends on $P_h$ through both arguments and does not simplify to a closed form for general $(\gamma_1,\gamma_2)$.  Resolving it would require Volterra inversion methods well beyond the scope of this thesis.

Model~$B$ is therefore left open at this point.  Its role in the hierarchy is structural: it identifies the precise obstruction (the Volterra equation for $P_y$) that the tractable special cases $B_2$ and $C_2$ circumvent, each by a different mechanism.  Model~$B_2$ sets $\gamma_2=0$, collapsing the PDE to an ODE whose boundary condition is pinned by an analyticity argument; Model~$C_2$ removes jockeying entirely and pins $P_y$ by a boundedness argument.  Both are analysed in the following sections.
